% !TEX root = ../thesis.tex

\chapter{Záver}

Cieľom tejto záverečnej práce bolo preskúmať možnosti detekcie toxického textu v prostredí webového prehliadača a navrhnúť, implementovať a vyhodnotiť riešenie vo forme rozšírenia \textit{Toxic Text Detector}. Práca sa zamerala na viacštítkovú klasifikáciu toxicity, používateľské rozhranie, ochranu súkromia, architektúru rozšírenia podľa Chrome Manifest V3 a porovnanie lokálneho a vzdialeného režimu inferencie. Výsledkom je funkčný prototyp hybridného systému, ktorý kombinuje spracovanie textu priamo v prehliadači s možnosťou vzdialenej inferencie.

V analytickej časti boli spracované teoretické východiská detekcie toxicity, relevantné datasety, modely, existujúce prístupy a technické aj etické problémy tejto oblasti. Analýza ukázala, že pri návrhu podobného systému nie je rozhodujúca iba presnosť modelu, ale aj spôsob nasadenia, latencia, transparentnosť výsledkov, ochrana súkromia a používateľská dôvera. Na základe týchto poznatkov bolo navrhnuté a implementované rozšírenie pozostávajúce z používateľského rozhrania v Reacte, service workera, content scriptu a vrstvy nastavení uchovávanej v úložisku prehliadača. Rozšírenie umožňuje analyzovať vybraný text alebo textové časti stránky, označiť potenciálne toxický obsah, vizuálne ho rozmazať a poskytnúť používateľovi základné vysvetlenie výsledku pomocou kategórie a skóre.

Dôležitou súčasťou práce bol aj návrh používateľského rozhrania. Rozhranie bolo vytvorené iteratívne, najskôr ako prototyp vo Figme a následne ako bežiaci HTML/React prototyp. Používateľské testovanie ukázalo, že zvolený koncept je zrozumiteľný a používatelia pozitívne vnímali najmä jednoduchosť ovládania, možnosť spätne zobraziť skrytý text a jasné signalizovanie stavu ochrany. Mechanizmus spätnej väzby používateľa pri nesprávnej klasifikácii bol zvažovaný v skorom návrhu, ale keďže išlo o voliteľné rozšírenie a jeho dopracovanie by si vyžadovalo ďalší návrh, ukladanie údajov a samostatné testovanie, nebol zahrnutý do finálnej implementácie.

V oblasti inferencie práca ukázala význam hybridného prístupu. Pôvodne bol ako lokálny variant overený predtrénovaný TF.js model toxicity, no jeho testovanie ukázalo obmedzenia z hľadiska rýchlosti, aktuálnosti a kvality výsledkov pri menej zastúpených kategóriách. Preto bol lokálny režim rozšírený o optimalizovaný TensorFlow.js GraphModel so samostatnými prahmi pre jednotlivé štítky a efektívnejším behom cez WASM backend. Lokálny režim sa ukázal ako rýchly a súkromie zachovávajúci základ, avšak s nižšou presnosťou pri náročnejších kategóriách. Vzdialený režim poskytol priestor pre presnejšiu detekciu a jednoduchšiu výmenu modelu, pričom pri lokálnom spustení API na výkonnom notebooku dosahoval odozvu porovnateľnú s lokálnym režimom. Vyhodnotenie zároveň potvrdilo, že pri detekcii toxicity je vhodné používať samostatné prahy pre jednotlivé kategórie namiesto jedného univerzálneho prahu.

Z pohľadu splnenia cieľov možno konštatovať, že hlavné ciele práce boli splnené. Bola spracovaná analytická časť, navrhnutá architektúra rozšírenia, vytvorené a otestované používateľské rozhranie, implementovaná funkčná verzia rozšírenia pre Chrome Manifest V3 a pripravený lokálny aj vzdialený režim inferencie. Práca zároveň ukázala obmedzenia riešenia, najmä pri detekcii implicitnej, kontextovej a sarkastickej toxicity, pri menej zastúpených kategóriách, pri viacjazyčných textoch a pri dôkladnom hodnotení spravodlivosti modelov. Ďalší vývoj by sa preto mohol zamerať na presnejšiu kalibráciu prahov, rozšírenie vzdialeného režimu o ďalšie modely, lepšie spracovanie rôznych jazykov, používateľskú spätnú väzbu a rozsiahlejšie testovanie na reálnych webových stránkach. Navrhnuté riešenie tak potvrdilo realizovateľnosť prehliadačovo integrovaného nástroja, ktorý vyvažuje presnosť detekcie, používateľskú kontrolu a ochranu súkromia.